% Valentino Vranic
% Metody inzinierskej prace 2012/13

\documentclass{beamer}

%\usetheme{Warsaw}
\usetheme{Antibes}
%\usetheme{JuanLesPins}
%\usetheme{Goettingen}

%\usecolortheme{seahorse}
%\usecolortheme{dolphin}
%\usecolortheme{rose}
% http://deic.uab.es/~iblanes/beamer_gallery/index_by_color.html
%\usecolortheme{beaver}
\usecolortheme{orchid}

%\useoutertheme[]{sidebar}

\setbeamercovered{transparent}

\usepackage[slovak]{babel}
\usepackage[T1]{fontenc}
\usepackage[utf8]{inputenc}
\usepackage{url}

\usepackage{listings}

\lstset{language=Python,basicstyle=\fontsize{8}{9.6}\selectfont,showstringspaces=false,columns=fullflexible,identifierstyle=\ttfamily,keywordstyle=\bfseries,showstringspaces=false,columns=fullflexible}
%\lstset{language=C,basicstyle=\fontsize{10.5}{12.6}\selectfont,identifierstyle=\ttfamily,keywordstyle=\bfseries,showstringspaces=false,columns=fixed}

\def\BibTeX{\textsc{Bib}\kern-.08em\TeX} 

\newcommand{\footcite}[1]{\footnote{\tiny #1}}
\newcommand{\umlet}{.5}
\newcommand{\emp}[1]{\textit{\alert{#1}}}
\newcommand{\kw}[1]{\mbox{\textbf{#1}}}
\newcommand{\id}[1]{\texttt{#1}}
\newcommand{\stl}{\guillemotleft}
\newcommand{\str}{\guillemotright}

\newcommand{\lsti}{\lstinline[basicstyle=\fontsize{10.5}{12.1}\selectfont]}

\newcommand{\ssection}[1]{
	\section{#1}
	\begin{frame}[fragile=singleslide]\frametitle{}
	\Huge #1
	\end{frame}
}

\newcommand{\ssectionn}[1]{
	\section*{#1}
	\begin{frame}[fragile=singleslide]\frametitle{}
	\Huge #1
	\end{frame}
}

\newenvironment{program}{\begin{beamercolorbox}[rounded=true,shadow=true]{block body}\vspace{-4mm}}{\vspace{-2mm}\end{beamercolorbox}}

\setbeamercolor{fvystup}{fg=white,bg=black}
\newenvironment{vystup}{\begin{beamercolorbox}[rounded=true,shadow=true]{fvystup}}{\end{beamercolorbox}}

\newenvironment{poznamka}{\begin{beamercolorbox}[rounded=true,shadow=false]{block body}}{\end{beamercolorbox}}

\setbeamertemplate{footline}[page number]
{
%\insertpagenumber
%\begin{beamercolorbox}{section in head/foot}
%\vskip2pt\insertnavigation{\paperwidth}\vskip2pt
%\end{beamercolorbox}%
}



\author{Štefan Kučerák}
%\url{www.fiit.stuba.sk/~vranic}, \url{vranic@fiit.stuba.sk}}
%{\tiny \url{www.fiit.stuba.sk/~vranic}, \url{vranic@fiit.stuba.sk}}
\institute{
	Ústav informatiky, informačných systémov a softvérového inžinierstva\\
	Fakulta informatiky a informačných technológií\\
	Slovenská technická univerzita v Bratislave}

\subtitle{\vspace{3mm} Metódy inžinierskej práce 2023/2024}

\title{Extrakcia informácii z webu pre analýzu dát pomocou Pythonu
}

\date{\footnotesize 6. november 2023}




\begin{document}

\begin{frame}[fragile=singleslide]
\titlepage
\end{frame}


\begin{frame}[fragile=singleslide]\frametitle{Úvod}

\begin{itemize}
\item automatizované zbieranie dát z web stránok
\item časová náročnosť
\end{itemize}

\end{frame}

\begin{frame}[fragile=singleslide]\frametitle{Prehľad}
\tableofcontents
\end{frame}


\section{Príklady reálneho využitia}
% príkaz \ssection by vytvoril zvláštný slajd s názvom časti - v krátkych prezentáciách to prekáža, lebo oberá o čas

\begin{frame}[fragile=singleslide]\frametitle{Príklady reálneho využitia}
\begin{itemize}
\item "Phishing Web Page Detection using web Scraping“
\item "Food Genie, Recipe Search Algorithm using Web Scraping"
\item "NEWSONE- AN AGGREGATION SYSTEM FOR NEWS USING WEB SCRAPING METHOD"
\end{itemize}
\end{frame}



\section{Knižnice}

\begin{frame}[fragile=singleslide]\frametitle{Knižnice}
\begin{itemize}
\item Requests\footcite{\url{https://pypi.org/project/requests/}}
\item Beautiful Soup\footcite{\url{https://pypi.org/project/beautifulsoup4/}}
\item Scrapy\footcite{\url{https://pypi.org/project/Scrapy/}}
\item Selenium\footcite{\url{https://pypi.org/project/selenium/}}

\end{itemize}
\end{frame}

\section{princíp}
\begin{frame}[fragile=singleslide]\frametitle{Princíp extrakcie dát}
\includegraphics[scale=0.75]{diag.png}
% pridajte vlastný obrázok a zrušte znák % pred príkazom \includegraphics vo formáte PDF prípadne PNG alebo JPG
% scale určuje veľkosť obrázku

{\tiny Nejaká poznámka k obrázku, možno zdroj\ldots}
\end{frame}

\section{príklad}
\begin{frame}[fragile=singleslide]\frametitle{Jednoduchý príklad}
\begin{poznamka}
Vypísanie názvu stránky pomocou knižnice Beautiful Soup
\end{poznamka}

\begin{itemize}
\item Program

\begin{program}
\begin{lstlisting}
import requests
from bs4 import BeautifulSoup

page = requests.get("https://www.google.com/")
soup = BeautifulSoup(page.content, "html.parser")

print(soup.title.string)
\end{lstlisting}
\end{program}

\item Výstup
\begin{vystup}
\begin{lstlisting}
Google
\end{lstlisting}
\end{vystup}

\end{itemize}
\end{frame}



\section*{Zhodnotenie a ďalšia práca}
% hviezdička zabezpečí, aby sa táto časť neocitla v prehľade prezentácie - každá prezentácia má zhodnotenie a prehľad by sa tým zbytočne zahlcoval

\begin{frame}[fragile=singleslide]\frametitle{Zhodnotenie a ďalšia práca}
\begin{itemize}
\item Každá prezentácia musí byť nejako uzavretá
\item Ale vždy je čo robiť ďalej\ldots{}
\end{itemize}
\end{frame}


\end{document}




Text \end{document} za príkazom \end{document} LaTeX ignoruje, takže tu môžete odkladať veci (aj celé slajdy), ktoré nechcete vymazať, lebo ich ešte možno budete potrebovať, avšak ich v danom momente nechcete mať v slajdoch.
